\documentclass[12pt]{article}
\usepackage[utf8]{inputenc}
\usepackage[T1]{fontenc}
\usepackage{hyperref}
\title{Riang Lai Language Resources for K-12 Curriculum Development}
\author{Compiled by Claude Sonnet 5.5 (claude-sonnet-5-5) for Andr\'as Kornai}
\date{}
\begin{document}
\maketitle

\section{Introduction}
Riang Lai (Yinchia; ISO 639-3 yin; Glottocode yinc1238) is an Austroasiatic language of northern Shan State, Myanmar, spoken in Namsan and Mantong Townships of the Palaung Self-Administered Zone. Joshua Project gives about 13,000 people in Myanmar, and a linguist's figure cited by the Max Planck Institute is about 12,000 speakers. It is closely related to Riang Lang (Yinnet, ISO ril), which should not be confused with it, and it is unrelated to the Riang of northeast India and Bangladesh. Sources note that many young people prefer Shan or Burmese. I found no educational materials.

\section{Primary Educational Resources}
None found.

\section{Linguistic and Reference Materials}
\begin{itemize}
\item Wikipedia, ``Riang language'': overview covering Riang Lai and Riang Lang. \url{https://en.wikipedia.org/wiki/Riang_language}
\item Max Planck Institute for Evolutionary Anthropology, Riang Lai numerals page. \url{https://lingweb.eva.mpg.de/channumerals/Riang-Lai.htm}
\end{itemize}

\section{Religious and Community Resources}
\begin{itemize}
\item Joshua Project profile of Riang Lai in Myanmar. \url{https://joshuaproject.net/people_groups/15091/BM}
\item Asia Harvest profile of Riang Lai (PDF). \url{https://www.operation.asiaharvest.org/MYANMAR-Riang-Lai.pdf}
\end{itemize}

\section{Digital and Audio Resources}
None found.

\section{Curriculum Development Context}
Almost nothing exists in the way of written materials: no orthography, dictionary, or primer turned up. Any curriculum would begin with community-led documentation and orthography development, probably alongside Shan or Burmese, which are the languages of wider communication. The related Riang Lang materials should not be assumed to apply.
\end{document}